\parthead{I}{A runway in a different direction}\label{part:1}
At Kai Tak, the radio beam that supplied lateral approach guidance did not point down the runway. This localiser pointed along $088^\circ$ magnetic. Runway 13 lay along $135^\circ$. Between the two directions stood the problem that made the airport famous: a large passenger aircraft had to make a substantial visual turn, descend, align over the pavement and arrive with enough room to stop \citep[§§1.8.1, 1.10]{cad}.

The setting made the geometry hard to ignore. Dense neighbourhoods crowded the landward side of the airport; hills rose beyond them; the runway extended into Kowloon Bay. From the city, the final turn looked improbably close. On a plan, it becomes a question about space and time. How much bank is needed to accumulate the change of direction? What height and speed can be carried through it? And, after the aircraft reaches the runway, how much of the landing remains?

\fig{01_place}{1}{Kai Tak and the city beside it in three separate Lands Department aerial frames, not a stitched panorama. CN4581 and CN4589 are dated 8 October 1993. A36143 is catalogued as 2 November 1993; its handwritten label appears to read 3 November, an unresolved discrepancy. Fine leaders locate visible features without inserting a flight path. The photographs retain their scan perspectives; no metric scale is inferred. Source: Lands Department, HKSAR Government; CSDI Portal and DATA.GOV.HK, DAP-L0. Government intellectual-property ownership retained; supplied as-is, without endorsement.}{fig:place}
\FloatBarrier

The two published bearings differ by $47^\circ$. They share a magnetic reference, so subtraction is legitimate. But the result describes directions in a horizontal plan. It does not say that the aircraft's wings must bank by 47 degrees. Bank is the inclination of the wings; a change of direction is the accumulated result of turning over time. Confusing the two makes the approach look like a much more violent manoeuvre than the number actually specifies.

The published geometry also separates two kinds of guidance. Kai Tak's Instrument Guidance System (IGS), introduced in 1974, supplied a localiser and a $3.1^\circ$ glidepath \citep{cadhistory,cad}. It brought an aircraft towards the visual segment. The pilot then had to make the right turn towards Runway 13, using the visible setting and the procedure. The instrument system had not bent its beam around the city.

\fig{02_bearings}{1}{The offset that defines the approach: $088^\circ$ magnetic for the IGS localiser and $135^\circ$ magnetic for Runway 13, a difference of $47^\circ$. Geographic context locates the problem; the bearing construction compares published directions, not a surveyed turn point. Magnetic bearing, ground track and air heading are distinct; the latter two separate in wind. Sources: CAD Report 1/95, §§1.8.1 and 1.10, Appendix 1; historical geographic sources identified in the figure provenance. Map fragment, landmark and runway symbol are schematic, not spatially registered.}{fig:bearings}

\begin{samepage}
The middle marker lay a reported 1.7 nautical miles from touchdown. If visual flight had not been achieved there, the procedure required a missed approach \citep[§1.8.1]{cad}. Visibility therefore belonged to the geometry of the task: the celebrated final turn was a visual manoeuvre, not an instrument-guided turn to be completed without sight of the landing environment. Checkerboard Hill supplied a visual reference in Kowloon, and the CAD's retrospective account places the right turn over Kowloon City \citep{cadlink}.
\par
\end{samepage}

A plan view can establish the directional mismatch. It cannot yet say whether the aeroplane has enough room. That depends on how quickly its velocity can turn and on where the turn begins and ends. Nor does alignment alone finish the approach. At the runway the aircraft must have a suitable height, descent rate and speed, all at the same time.

The final distance has a specific historical definition. CAD Report 1/95 gives a runway length of 3,331.5 m, a width of 61 m and a landing distance available of 2,786 m \citep[§1.10]{cad}. The last quantity is the one used in the threshold-relative stopping calculations below. The longer lengths reported in general airport histories describe a different quantity or historical context; Appendix~\ref{app:sources} retains those distinctions. A landing cannot spend pavement that its distance convention has already excluded.

By the time Kai Tak closed in July 1998 and Hong Kong's airport moved to Chek Lap Kok, this arrangement had become one of aviation's most recognizable scenes \citep{cadlink}. Its mechanics begin with a less dramatic image: the lift vector of a steadily banked aircraft.
